% Phase Coherence as an Admissibility Constraint
% Companion paper to the Cosmochrony framework and the O-series
% spectral admissibility sub-programme.

\documentclass[11pt,a4paper]{article}

\usepackage{amsmath,amssymb,amsthm}
\usepackage{enumerate}
\usepackage{hyperref}
\usepackage{doi}
\usepackage{graphicx}

\newtheorem{theorem}{Theorem}[section]
\newtheorem{lemma}[theorem]{Lemma}
\newtheorem{corollary}[theorem]{Corollary}
\newtheorem{proposition}[theorem]{Proposition}
\newtheorem{remark}[theorem]{Remark}
\newtheorem{definition}[theorem]{Definition}

\title{From Admissibility to Quantum Structure:\\
Phase Coherence and Correlations from Non-Injective Projection}

\author{J\'er\^ome Beau\\
\small Independent Researcher, Paris, France\\
\small \texttt{jerome.beau@cosmochrony.org}}

\date{2026}

\begin{document}

  \maketitle

  \begin{abstract}
    We establish that quantum-mechanical phase coherence is not an independent
    postulate but a structural consequence of projective admissibility within
    the Cosmochrony framework.
    The central result (Theorem~\ref{thm:admissibility-phase-coherence}) shows
    that any admissible transition preserves the Born--Infeld indiscernibility of
    conjugate Weil blocks $\rho_c$ and $\rho_{q-c}$ on the Heisenberg Cayley graph
    $\mathrm{Cay}(\mathrm{Heis}_3(\mathbb{Z}/q\mathbb{Z}), S_q)$, and thereby
    maintains the metaplectic phase coherence of the admissible fibre.
    The proof rests on three independently established results: the minimal parity
    fibre structure of the non-injective projection (O18), the dynamical identity of
    conjugate Weil blocks (O17), and the discrete admissibility filter of projection
    locking (O22).
    A constructive separating functional $R_n^{(c)}$ is exhibited and proved
    BI-admissible.
    Numerical validation for $q=29$ across four conjugate pairs confirms
    that $\mathrm{rank}\,W_n^{(c)} = 0$ exactly throughout the admissible regime,
    establishing structural indistinguishability of conjugate blocks as a
    representation-theoretic fact rather than a dynamical coincidence.
    As a structural consequence, the admissible fibre carries complex amplitude
    structure with interference, constituting the natural precursor of an effective
    Hilbert sector — without quantum mechanics as a postulate.
    The singlet correlator $E(\hat{a},\hat{b}) = -\hat{a}\cdot\hat{b}$ is then derived
    from admissibility, the parity involution (O18), and the isotropic SU(2) structure (O23),
    without invoking any quantum postulate.
    As a consequence, the Tsirelson bound $|S_{\mathrm{CHSH}}| \leq 2\sqrt{2}$ follows
    unconditionally from the geometry of $\mathrm{Im}\,\mathbb{H}$ under admissible projection.
    Finally, the Born probability rule $P(a=+1\,|\,\hat{a}) = |\langle+\hat{a}|\psi\rangle|^2$
    is derived as the unique probability assignment compatible with the coherent fibre structure,
    the SU(2) observable algebra, and the two-point correlator law.
    No quantum postulate is invoked at any step of the derivation.
  \end{abstract}

  \paragraph{Keywords.}
    Non-injective projection; Quantum foundations; Emergence; Phase coherence; Born–Infeld structure; Weil
    representation; SU(2) symmetry; Fiber structure; Admissibility; Gram–Schmidt dynamics

    \tableofcontents

    \clearpage

  \section{Introduction}
  \label{sec:introduction}

  The structural origin of quantum-mechanical phase coherence is one of the deepest
  open questions in the foundations of physics.
  Standard formulations of quantum mechanics postulate the Hilbert space structure,
  complex amplitudes, and the superposition principle as axioms.
  The present paper shows that, within the projective admissibility framework of
  Cosmochrony~\cite{Cosmochrony}, phase coherence is not postulated but derived:
  it is a necessary consequence of the admissibility constraint that governs
  which transitions between observable states are physically realised.

  \subsection*{Context and position in the programme}

    The Cosmochrony framework~\cite{Cosmochrony} derives physical structure from
    a single primitive: the non-injective projection between successive observable
    states.
    Non-injectivity is not a modelling choice but a structural necessity — any
    projection compatible with genuine emergence must be non-injective, as
    established in~\cite{ENI}.
    The Born--Infeld structure is selected as the unique admissible completion
    of the effective action compatible with bounded projective flux~\cite{BornInfeld}.

    Within this framework, the spectral admissibility sub-programme (O-series)
    has established a concrete chain of results on the Heisenberg Cayley graph
    $\mathrm{Cay}(\mathrm{Heis}_3(\mathbb{Z}/q\mathbb{Z}), S_q)$:
    the minimal parity fibre $\{c, q-c\}$ forced by Born--Infeld evenness (O18),
    the discrete admissibility filter of projection locking (O22),
    and the derivation of exactly three stable neutral directions from quaternionic
    minimality (O23).

    The present paper establishes the bridge between this spectral programme and
    quantum mechanics.
    It answers the question: why does the admissible fibre carry complex amplitudes
    and support interference, rather than being a mere set of classical alternatives?

    The present paper sits at the intersection of three lines of work:
    the structural failure of Bell factorisation from non-injective projection~\cite{Bell};
    the spectral admissibility sub-programme (O-series), which establishes concrete
    results on the Heisenberg Cayley graph including the parity fibre (O18),
    projection locking (O22), and quaternionic minimality (O23);
    and the new foundational model of admissible non-injective transitions~\cite{Cosmochrony}.
    All three are necessary inputs.
    From the Bell paper, the framework inherits the structural argument that
    non-injectivity makes Bell factorisation inapplicable.
    From the O-series, it inherits the concrete spectral tools that make the
    admissibility constraint explicit.
    From the foundational model, it inherits the physical principle that
    admissibility means non-premature selection.
    The present paper connects these three lines: it shows that the admissibility
    constraint, made concrete by the O-series, forces the fibre to carry the
    coherent structure that makes quantum predictions possible.

  \subsection*{Main results}

    The central results of this paper are the following.

    \begin{enumerate}[(i)]
\item \textbf{Admissibility implies phase coherence}
(Theorem~\ref{thm:admissibility-phase-coherence}).
Any admissible transition preserves the Born--Infeld indiscernibility of
conjugate Weil blocks $\rho_c$ and $\rho_{q-c}$, and thereby maintains their
metaplectic phase coherence.
Phase coherence is not preserved because it is postulated — it is preserved
because destroying it would constitute a premature selection among admissible
directions, violating the admissibility constraint.

\item \textbf{Observable signature}
(Corollary~\ref{cor:gram-schmidt-rank}).
Any coherence-destroying transition induces a rank inflation in the combined
Gram--Schmidt span of $\{\rho_c, \rho_{q-c}\}$, incompatible with the admissible
rank growth of O22.
This provides a falsifiable signature detectable in the spectral pipeline.

\item \textbf{Singlet correlator from admissibility}
(Theorem~\ref{thm:singlet-correlator}).
The correlator $E(\hat{a},\hat{b}) = -\hat{a}\cdot\hat{b}$ is derived from
four structural inputs: the linear observable structure of $\mathfrak{su}(2)$
(O23), the bilinear evaluation induced by coherent fibre amplitudes (Q1),
rotation invariance from the isotropic condition $M \propto I$ (O23),
and normalisation from the parity involution (O18).
No quantum postulate is invoked.

\item \textbf{Tsirelson bound — unconditional}
(Corollary~\ref{cor:tsirelson-unconditional}).
Since the singlet correlator is now derived rather than assumed,
$|S_{\mathrm{CHSH}}| \leq 2\sqrt{2}$ follows unconditionally from the
geometry of $\mathrm{Im}\,\mathbb{H}$ under admissible projection.

\item \textbf{Born rule from admissibility}
(Theorem~\ref{thm:born-rule}).
The Born probability rule $P(a=+1\,|\,\hat{a}) = |\langle+\hat{a}|\psi\rangle|^2$
is derived as the unique probability assignment consistent with:
positivity and normalisation, linearity of mean values (from Q1 coherence),
and compatibility with the derived correlator.
No Gleason theorem~\cite{Gleason1957} is needed: uniqueness follows from the structural
constraints already established.
    \end{enumerate}

  \subsection*{What this paper does not claim}

    This paper does not derive the full Hilbert space structure for arbitrary
    observables, multipartite systems, or continuous spectra.
    Within the SU(2) sector, however, the derivation is structurally complete:
    admissibility forces coherence (Theorem~\ref{thm:admissibility-phase-coherence}),
    coherence combined with SU(2) geometry forces the singlet correlator
    (Theorem~\ref{thm:singlet-correlator}), the singlet correlator implies the
    Tsirelson bound (Corollary~\ref{cor:tsirelson-unconditional}), and the
    correlator combined with coherence uniquely forces the Born rule
    (Theorem~\ref{thm:born-rule}).
    The extension of these results beyond the SU(2) parity sector remains open.

  \subsection*{Organisation}

    Section~\ref{sec:admissibility-phase-coherence} establishes the main results:
    the key lemma linking metaplectic phase coherence to Born--Infeld
    indiscernibility (Section~\ref{subsec:q1-lemma}),
    the main theorem (Section~\ref{subsec:q1-theorem}),
    the Gram--Schmidt rank corollary (Section~\ref{subsec:q1-corollary}),
    the consequences for quantum structure (Section~\ref{subsec:q1-consequences}),
    the observable signature of fibre separation
    (Section~\ref{subsec:observable-signature}),
    the main theorem (Section~\ref{subsec:q1-theorem}),
    the Gram--Schmidt rank corollary (Section~\ref{subsec:q1-corollary}),
    the consequences for quantum structure (Section~\ref{subsec:q1-consequences}),
    the conditional Tsirelson bound (Section~\ref{subsec:tsirelson}),
    the derivation of the singlet correlator (Section~\ref{subsec:singlet-correlator}),
    the unconditional Tsirelson bound (Corollary~\ref{cor:tsirelson-unconditional}),
    and the Born rule (Section~\ref{subsec:born-rule}).

% -----------------------------------------------------------------------
% BODY (imported from cosmochrony-Q1-theorem.tex)
% -----------------------------------------------------------------------

    \input{Q1-theorem}

% -----------------------------------------------------------------------

  \section{Conclusion}
  \label{sec:conclusion}

  This paper has established that phase coherence in the admissible fibre is a
  structural consequence of the admissibility constraint, not an independent
  quantum postulate.
  The argument rests on four independently established pillars:
  the Born--Infeld minimal parity fibre (O18),
  the dynamical identity of conjugate Weil blocks (O17),
  the metaplectic central-phase coherence structure (O12, O19),
  and the projection locking theorem (O22).

  The central result, Theorem~\ref{thm:admissibility-phase-coherence}, closes
  the structural gap identified as Q1: any admissible transition preserves
  Born--Infeld indiscernibility of conjugate Weil blocks, and thereby forces the
  fibre to carry coherent complex amplitudes.
  The Gram--Schmidt rank corollary provides a concrete, falsifiable signature
  of any violation.
  Numerical validation for $q=29$ across four conjugate pairs confirms
  $\mathrm{rank}\,W_n^{(c)} = 0$ exactly throughout the admissible regime
  (Section~\ref{subsec:observable-signature}): structural indistinguishability
  is universal across all conjugate pairs and is a representation-theoretic
  identity, not a dynamical coincidence.

  Theorem~\ref{thm:singlet-correlator} then derives the singlet correlator
  $E(\hat{a},\hat{b}) = -\hat{a}\cdot\hat{b}$ from four structural inputs:
  linearity of $\mathfrak{su}(2)$ observables (O23),
  bilinearity of evaluation from coherent fibre amplitudes (Q1),
  rotation invariance from isotropic saturation (O23),
  and normalisation from the parity involution (O18).
  No quantum postulate is invoked at any step.

  As a consequence, Corollary~\ref{cor:tsirelson-unconditional} establishes the
  Tsirelson bound $|S_{\mathrm{CHSH}}| \leq 2\sqrt{2}$ unconditionally:
  it follows from the geometry of $\mathrm{Im}\,\mathbb{H}$ under admissible
  projection, not from any quantum axiom.

  Theorem~\ref{thm:born-rule} then completes the structural derivation:
  the Born rule $P(a=+1\,|\,\hat{a}) = |\langle+\hat{a}|\psi\rangle|^2$
  is the unique probability assignment compatible with the coherent fibre structure,
  the SU(2) observable algebra, and the derived correlator.
  No Gleason theorem~\cite{Gleason1957} is needed.
  For the SU(2) sector, quantum mechanics is not postulated — it is forced
  by the admissibility structure of the non-injective projection.

  \subsection*{Open directions}

    Three directions remain open after this paper.

    \paragraph{Extension of numerical validation to larger primes.}
      The structural indistinguishability $\mathrm{rank}\,W_n^{(c)} = 0$
      has been confirmed numerically for $q=29$.
      Systematic verification across the full prime range of O25
      ($q \in \{29, 61, 101, 151, 211, 307, 401\}$) would provide a comprehensive
      empirical confirmation and allow tracking of the boundary transition
      (rank-1 residual at saturation) as a function of $q$.

    \paragraph{Extension of the Born rule beyond SU(2).}
      Theorem~\ref{thm:born-rule} derives the Born rule for spin-$\tfrac{1}{2}$
      observables in the SU(2) sector.
      The extension to arbitrary observables, multipartite systems, and
      continuous spectra requires extending the fibre structure beyond the
      two-element parity case of O18.
      This generalisation is structurally motivated and is the primary remaining
      open direction of the quantum bridge programme.

    \paragraph{BI-admissibility of the separating functional.}
      The converse direction of Lemma~\ref{lem:coherence-indiscernibility} establishes
      the existence of a Gram--Schmidt-based functional distinguishing the two blocks
      in the incoherent regime.
      A fully rigorous construction of this functional, together with an explicit
      verification of its Born--Infeld admissibility, would strengthen the lemma from
      an existence result to a constructive one.

  \subsection*{Position in the programme}

  The result of this paper occupies a specific position in the Cosmochrony chain:
  \begin{center}
    \begin{tabular}{ccc}
      \cite{Bell}                           &            & \\
      (NI $\Rightarrow$ Bell fails)         & $\searrow$ &                                         \\
      &            & \textbf{this paper}                     \\
      O-series (O17, O18, O19, O22, O23)    & $\nearrow$ & (admissibility $\Rightarrow$ coherence)\\
      (spectral concrete results)           &            &                                         \\
      $\uparrow$                            &            &                                         \\
      \cite{Cosmochrony} foundational model &            &
    \end{tabular}
  \end{center}
  It provides the missing structural link explaining why the fibre is not merely
  a set of classical alternatives, but carries the coherent amplitude structure
  that makes quantum predictions possible.

  \section{Comparison with standard quantum mechanics}
  \label{sec:comparison-qm}

  The results of this paper derive quantum-mechanical structure without invoking
  any of the standard postulates of quantum theory.
  Table~\ref{tab:comparison} makes the comparison explicit.

  \begin{table}[h]
    \centering
    \begin{tabular}{lll}
      \hline
      \textbf{Element}        & \textbf{Standard QM}         & \textbf{This paper}                                   \\
      \hline
      Hilbert space           & Postulated                   & Precursor derived (SU(2) sector)                      \\
      Complex amplitudes      & Postulated                   & Forced by admissibility (Q1)                          \\
      Superposition principle & Postulated                   & Follows from fibre coherence                          \\
      Measurement outcomes    & Postulated ($\pm 1$)         & Spectrum of $\hat\sigma\cdot\hat{a}$                  \\
      Correlator law          & Postulated (QM rule)         & Derived (Theorem~\ref{thm:singlet-correlator})        \\
      Tsirelson bound         & Theorem~\cite{Tsirelson1980} & Derived (Corollary~\ref{cor:tsirelson-unconditional}) \\
      Born rule               & Postulated~\cite{Born1926}   & Derived (Theorem~\ref{thm:born-rule})                 \\
      Non-locality            & Assumed or interpreted       & Structural (non-injective projection)                 \\
      \hline
    \end{tabular}
    \caption{Comparison between standard quantum mechanics and the present
    admissibility-based derivation. Elements marked ``derived'' follow from
    the admissibility constraint and the O-series structural results,
      without quantum postulates as input.}
    \label{tab:comparison}
  \end{table}

  Three structural differences from standard quantum mechanics are worth
  making explicit.

  \paragraph{No Hilbert space postulate.}
    Standard quantum mechanics postulates the Hilbert space $\mathcal{H}$,
    complex amplitudes, and the superposition principle.
    Here, the admissible fibre $F_n = \Pi_n^{-1}(O_n)$ carries complex
    amplitudes because any admissible transition must preserve the
    Born--Infeld indiscernibility of conjugate blocks
    (Theorem~\ref{thm:admissibility-phase-coherence}).
    The Hilbert space structure is a consequence, not a starting point.

  \paragraph{No local hidden-variable completion.}
    The present framework is not a hidden-variable model~\cite{Bell1964}.
    The underlying configuration is not a set of pre-determined outcome values —
    it is an admissible non-injective transition whose fibre structure prevents
    factorization.
    Bell factorizability fails not because of fine-tuned correlations between
    settings and hidden states, but because the fibre $F_n$ is non-decomposable
    (see~\cite{Bell}).
    Statistical independence of measurement settings is preserved.

  \paragraph{Quantum structure is derived, not compatible.}
    Previous structural approaches (e.g.\ ontological models~\cite{Harrigan2010},
    epistemic interpretations~\cite{Spekkens2007}) typically aim to show that quantum mechanics is
    \emph{compatible with} some underlying structure.
    Here the aim is stronger: quantum-mechanical correlations are
    \emph{forced} by the admissibility constraint,
    in the spirit of reconstruction programmes~\cite{Chiribella2011}
    but from a structurally different starting point.
    Any violation of the Born rule, the singlet correlator, or the
    Tsirelson bound would require a transition that is not admissible —
    i.e.\ one that constitutes a premature selection among open projective
    directions, violating Axiom~A3.

  \section{Falsifiability}
  \label{sec:falsifiability}

  The results of this paper generate falsifiable predictions at two levels:
  within the existing spectral pipeline, and at the level of quantum
  correlations.

  \paragraph{Pipeline-level prediction.}
    The central prediction is:
    \begin{equation}
      \mathrm{rank}\,W_n^{(c)} = 0
      \quad \text{for all conjugate pairs } (c, q-c)
      \text{ and all depths } n \in [n_0, n_1].
    \end{equation}
    This follows from the representation-theoretic identity
    $\rho_{q-c} = \overline{\rho_c}$, which implies identical Gram--Schmidt
    spans.
    Any deviation — i.e.\ $\mathrm{rank}\,W_n^{(c)} > 0$ in the admissible
    window — would falsify the structural indistinguishability of conjugate
    blocks, and thereby invalidate the foundation of
    Theorem~\ref{thm:admissibility-phase-coherence}.

    Numerical validation for $q = 29$ across four conjugate pairs confirms
    $\mathrm{rank}\,W_n^{(c)} = 0$ exactly throughout $[n_0, n_1]$
    (Section~\ref{subsec:observable-signature}).
    Extension to larger primes $q \in \{61, 101, 151, 211, 307, 401\}$
    is an immediate empirical target.

  \paragraph{Correlator-level prediction.}
    Theorem~\ref{thm:singlet-correlator} predicts:
    \begin{equation}
      E(\hat{a}, \hat{b}) = -\hat{a} \cdot \hat{b}
    \end{equation}
    for any two-point measurement on a shared admissible fibre.
    A measured correlator deviating from this law — in particular,
    a CHSH parameter exceeding $2\sqrt{2}$, or a correlator following
    the classical sign-projection law $1 - 2\theta/\pi$ — would falsify
    the SU(2) coherent structure of the admissible fibre.

  \paragraph{Born rule prediction.}
    Theorem~\ref{thm:born-rule} predicts:
    \begin{equation}
      P(a = +1\,|\,\hat{a}) = |\langle +\hat{a}\,|\,\psi\rangle|^2.
    \end{equation}
    Any probability assignment deviating from this law while maintaining
    coherent amplitude structure and the derived correlator is ruled out
    by the uniqueness argument of Step~3.
    In particular, there is no free parameter: the exponent $2$ in the
    Born rule is not a convention but a structural necessity.

  \paragraph{Summary.}
    The framework generates three independent falsifiable predictions —
    GS rank, correlator law, and Born weights — each testable at a
    different level of the pipeline.
    If any one fails, the admissibility-based derivation of quantum
    mechanics is falsified.
    If all three hold, they constitute joint evidence that quantum structure
    is a structural consequence of projective admissibility.
    The correlator and Born predictions are consistent with all existing
    experimental tests of quantum mechanics~\cite{Aspect1982}.

  \section*{Acknowledgements}

  The author acknowledges the use of large language models as a supportive tool
  for refining language, structure, and internal consistency during the development
  of this manuscript.
  All conceptual contributions, theoretical choices, and interpretations remain
  the sole responsibility of the author.

  \clearpage
  \begin{thebibliography}{99}

    \bibitem{Born1926}
    M.~Born,
    ``Zur Quantenmechanik der Sto\ss{}vorg\"ange,''
    \textit{Zeitschrift f\"ur Physik} \textbf{37}, 863--867 (1926).
    \doi{10.1007/BF01397477}.

    \bibitem{Tsirelson1980}
    B.~S.~Tsirelson,
    ``Quantum generalizations of Bell's inequality,''
    \textit{Letters in Mathematical Physics} \textbf{4}, 93--100 (1980).
    \doi{10.1007/BF00417500}.

    \bibitem{Gleason1957}
    A.~M.~Gleason,
    ``Measures on the closed subspaces of a Hilbert space,''
    \textit{Journal of Mathematics and Mechanics} \textbf{6}, 885--893 (1957).
    \doi{10.1512/iumj.1957.6.56050}.

    \bibitem{Bell1964}
    J.~S.~Bell,
    ``On the Einstein Podolsky Rosen paradox,''
    \textit{Physics Physique Fizika} \textbf{1}, 195--200 (1964).
    \doi{10.1103/PhysicsPhysiqueFizika.1.195}.

    \bibitem{Aspect1982}
    A.~Aspect, P.~Grangier, and G.~Roger,
    ``Experimental realization of Einstein-Podolsky-Rosen-Bohm
    \textit{Gedankenexperiment}: a new violation of Bell's inequalities,''
    \textit{Physical Review Letters} \textbf{49}, 91--94 (1982).
    \doi{10.1103/PhysRevLett.49.91}.

    \bibitem{Spekkens2007}
    R.~W.~Spekkens,
    ``Evidence for the epistemic view of quantum states: a toy theory,''
    \textit{Physical Review A} \textbf{75}, 032110 (2007).
    \doi{10.1103/PhysRevA.75.032110}.

    \bibitem{Harrigan2010}
    N.~Harrigan and R.~W.~Spekkens,
    ``Einstein, incompleteness, and the epistemic view of quantum states,''
    \textit{Foundations of Physics} \textbf{40}, 125--157 (2010).
    \doi{10.1007/s10701-009-9347-0}.

    \bibitem{Chiribella2011}
    G.~Chiribella, G.~M.~D'Ariano, and P.~Perinotti,
    ``Informational derivation of quantum theory,''
    \textit{Physical Review A} \textbf{84}, 012311 (2011).
    \doi{10.1103/PhysRevA.84.012311}.

    \bibitem{Cosmochrony}
    J.~Beau,
    ``Cosmochrony: A Pre-Geometric Relational Framework for Emergent Physics,''
    white paper, Zenodo (2026).
    \doi{10.5281/zenodo.17957509}.

    \bibitem{ENI}
    J.~Beau,
    ``Non-Injectivity as a Structural Necessity of Genuine Emergence:
    A No-Go Theorem for Faithful Emergent Descriptions,''
    Cosmochrony companion paper~L, Zenodo (2026).
    \doi{10.5281/zenodo.19391746}.

    \bibitem{BornInfeld}
    J.~Beau,
    ``Bounded Relaxation and the Dynamical Selection of Spacetime Geometry,''
    Cosmochrony companion paper~C, Zenodo (2026).
    \doi{10.5281/zenodo.18407505}.

    \bibitem{Bell}
    J.~Beau,
    ``Non-Injective Projections and the Structural Failure of Bell Factorization,''
    Cosmochrony companion paper~B, Zenodo (2026).
    \doi{10.5281/zenodo.18371173}.

    \bibitem{O12}
    J.~Beau,
    ``Exact Weil-Block Projective Capacity on Heisenberg Graphs,''
    O12 spectral admissibility sub-programme paper, Zenodo (2026).
    \doi{10.5281/zenodo.19194798}.

    \bibitem{O17}
    J.~Beau,
    ``Fibre-Level Admissibility from Conjugate Weil Blocks:
    Structural Derivation of the Pair Observable,''
    O17 spectral admissibility sub-programme paper, Zenodo (2026).
    \doi{10.5281/zenodo.19299611}.

    \bibitem{O18}
    J.~Beau,
    ``Minimal Fibre Structure of the Non-Injective Projection
    from Born--Infeld Indiscernibility: Derivation of the Parity Involution,''
    O18 spectral admissibility sub-programme paper, Zenodo (2026).
    \doi{10.5281/zenodo.19321413}.

    \bibitem{O19}
    J.~Beau,
    ``Canonical Pair Observables in Weil Blocks: Structure of the Residual
    Amplitude Factor and Pipeline-Independent Formulation,''
    O19 spectral admissibility sub-programme paper, Zenodo (2026).
    \doi{10.5281/zenodo.19323720}.

    \bibitem{O22}
    J.~Beau,
    ``Shell-Alignment from Projection Locking: A Discrete Admissibility Theorem
    under Born--Infeld Saturation,''
    O22 spectral admissibility sub-programme paper, Zenodo (2026).
    \doi{10.5281/zenodo.19367375}.

    \bibitem{O23}
    J.~Beau,
    ``Three Stable Directions from Quaternionic Minimality: Derivation of
    $\Sigma_c(n_3) = 3$ from Born--Infeld Fibre Admissibility,''
    O23 spectral admissibility sub-programme paper, Zenodo (2026).
    \doi{10.5281/zenodo.19375136}.

  \end{thebibliography}

\end{document}
